\section{Introduction}
\label{sec:introduction}

Fifth-generation wireless systems use multiple frequency ranges, and planar microstrip antennas are
widely studied for compact sub-6-GHz radio-frequency front ends. Rectangular microstrip patch
antennas are attractive because they are low profile, planar, mechanically simple, and compatible
with printed-circuit fabrication. Their principal limitations include relatively narrow impedance
bandwidth and sensitivity to substrate properties, feed geometry, conductor loss, and fabrication
tolerances.

This report treats the antenna as a complete engineering design problem. The CST model is not used
as a black box; instead, the geometry is initialized from microstrip antenna theory, tuned in
full-wave simulation, and interpreted using performance metrics such as \(S_{11}\), radiation
pattern, directivity, radiation efficiency, and total efficiency. The tuned patch is then used as the
starting point for a two-element array. The array analysis connects full-wave simulation to phased-array
theory through element spacing, array factor, mutual coupling, and beam steering under controlled
relative excitation phase.
